\documentclass{article}
\usepackage[T1]{fontenc}
\usepackage{lmodern}
\usepackage{graphicx}
\usepackage{amsmath,amssymb}
\usepackage[a4paper,margin=2cm]{geometry}
\usepackage[absolute,overlay]{textpos}
\setlength{\TPHorizModule}{1mm}
\setlength{\TPVertModule}{1mm}
\usepackage[indonesian]{babel}
\usepackage{hyperref}
\setlength{\emergencystretch}{2em}
% unit: Halaman 11

\begin{document}

% === Halaman 11 (python/native) ===

11

Prosedur Enkripsi 

1.

Preprocessing: Jika plainteks m berukuran besar, pecah plainteks 

menjadi blok-blok m1, m2, …, (nilai setiap blok harus berada di dalam 

selang [0, p – 1].

2.

Pilih bilangan acak k, yang dalam hal ini 1  k  p – 2. 

3.

Setiap blok m dienkripsi dengan rumus



a = gk mod p  





(2)



b = ykm mod p 





(3)


Pasangan (a, b) adalah cipherteks untuk blok pesan m. Jadi, ukuran 

cipherteks adalah dua kali ukuran plainteksnya.

\end{document}
